\chapter{\texorpdfstring{\(L^p\)}{Lp} 空间中的逼近}

\begin{solution}{simple-dense-complex}
写 \(f=u+iv\)，其中 \(u,v\) 为实值且
\(|u|,|v|\le|f|\)，故都属于 \(L^p\)。分别取实简单函数
\(u_n\to u,v_n\to v\) 于 \(L^p\)，令
\(\varphi_n=u_n+iv_n\)。复简单函数只取有限多个复数值；由
\(|a+ib|\le|a|+|b|\) 及实数的幂次估计，
\[
 \|\varphi_n-f\|_p^p
 \le C_p\bigl(\|u_n-u\|_p^p+\|v_n-v\|_p^p\bigr)\to0,
\]
其中 \(C_p=1\)（\(0<p\le1\)），\(C_p=2^{p-1}\)（\(p\ge1\)）。
因此不必预先建立复标量的 Minkowski 不等式。
\end{solution}

\begin{solution}{finite-support-simple-dense}
先设 \(f\ge0\)，并沿用非负简单逼近中的二进构造
\[
 \varphi_n(x)=
 \begin{cases}
  2^{-n}\lfloor2^nf(x)\rfloor,&f(x)<2^n,\\
  2^n,&f(x)\ge2^n.
 \end{cases}
\]
于是 \(0\le\varphi_n\uparrow f\)，而
\[
 \{\varphi_n\ne0\}\subseteq\{f\ge2^{-n}\},
 \qquad
 \mu(\{f\ge2^{-n}\})
 \le2^{np}\|f\|_p^p<\infty.
\]
又有 \(|f-\varphi_n|^p\le f^p\) 且逐点趋零，控制收敛定理给出
\(\|f-\varphi_n\|_p\to0\)。一般实值函数分别处理正部、负部；复值函数再
分别处理实部、虚部。整个论证没有使用 \(X\) 的 \(\sigma\)-有限分解。
\end{solution}

\begin{solution}{compact-open-distance-function}
引理在 \(K=\varnothing\) 时已经取 \(g=0\)，不使用距离比值。
以下 \(K\ne\varnothing\)；\(U\) 有界，所以 \(U^c\ne\varnothing\)。
到闭集的距离函数是 \(1\)-Lipschitz，故分子、分母连续。若分母为零，则
某点同时属于 \(K\) 与 \(U^c\)，与 \(K\subset U\) 矛盾。因此商连续。
在 \(K\) 上 \(d(x,K)=0\)，故 \(g=1\)；在 \(U^c\) 上
\(d(x,U^c)=0\)，故 \(g=0\)。又 \(0\le g\le1\)，且其非零集包含于
\(U\)，所以支撑包含于有界闭集 \(\overline U\)，为紧集。
\end{solution}

\begin{solution}{cc-density-error-budget}
删去零系数项，并记剩余项数为 \(m\)。若 \(m=0\)，取 \(g=0\) 即可。
以下 \(m\ge1\)。若 \(p\ge1\)，对每个 \(j\) 由
引理取 \(g_j\in C_c\)，使
\[
 \|\mathbf1_{E_j}-g_j\|_p
 <\frac{\varepsilon}{m(1+|a_j|)}
\]
并令 \(g=\sum_ja_jg_j\)。Minkowski 给
\[
 \|\varphi-g\|_p
 \le\sum_j|a_j|\|\mathbf1_{E_j}-g_j\|_p
 <\varepsilon.
\]
若 \(0<p<1\)，改取
\[
 \|\mathbf1_{E_j}-g_j\|_p
 <\frac{\varepsilon}{m^{1/p}(1+|a_j|)}.
\]
利用 \(|\sum_jb_j|^p\le\sum_j|b_j|^p\)，可得
\[
 \|\varphi-g\|_p^p
 \le\sum_j|a_j|^p\|\mathbf1_{E_j}-g_j\|_p^p
 <\varepsilon^p.
\]
\(g\) 是有限个连续紧支撑函数之和，仍属于 \(C_c\)。
\end{solution}

\begin{solution}{lp-separable-rational-grid}
令 \(\mathcal S\) 为所有有限和
\[
 \sum_{j=1}^{m}q_j\mathbf1_{R_j},
\]
其中 \(q_j\in\mathbb Q\)，\(R_j\) 是有理端点半开有界矩形。所有有理数
和有理矩形构成可数集；可数集的有限笛卡儿积及其可数并仍可数，故
\(\mathcal S\) 至多可数，且包含无限多个不同常值倍数的单个矩形指示函数，
所以可数无限。对 \(g\in C_c\)，取 \(Q\) 使支撑包含于其内部。
给定 \(\delta>0\)，网格足够细时，每格的振幅小于 \(\delta/2\)；
选样点值的有理近似，误差也小于 \(\delta/2\)。所得 \(s\in\mathcal S\)
在 \(Q\) 外取零，且 \(\|s-g\|_p^p\le m(Q)\delta^p\)。外层面零测。
再用 \(C_c\) 稠密性即得。若讨论复值空间，应将系数换成
\(\mathbb Q+i\mathbb Q\)，它仍可数无限。
\end{solution}

\begin{solution}{linfty-nonseparable-balls}
若 \(D\) 是可数稠密集，则每个
\(\mathbf1_{(0,t)}\) 可选 \(g_t\in D\)，使
\(\|g_t-\mathbf1_{(0,t)}\|_\infty<1/3\)。若 \(s\ne t\) 却有
\(g_s=g_t\)，三角不等式将给
\[
 1=\|\mathbf1_{(0,s)}-\mathbf1_{(0,t)}\|_\infty<2/3,
\]
矛盾。因此 \(t\mapsto g_t\) 是从不可数集到可数集的单射，不可能。
\end{solution}

\begin{solution}{translation-isometry}
平移不变性先给出指示函数积分不变，再由简单逼近和单调收敛推广到非负
可测函数。因而对 \(0<p<\infty\)，
\[
 \|\tau_hf\|_p^p
 =\int|f(x-h)|^p\,dx
 =\int|f(y)|^p\,dy.
\]
\(p=\infty\) 时零测集在平移下仍零测，同样成立。一般测度若不平移不变，
平移前后的积分未必相等。例如对 \(d\mu=e^{-x^2}dx\)，取
\(f=\mathbf1_{(0,1)}\)、\(h=1\)。则
\[
 \|\tau_1f\|_{L^p(\mu)}^p=\int_1^2 e^{-x^2}\,dx
 <\int_0^1 e^{-x^2}\,dx=\|f\|_{L^p(\mu)}^p.
\]
这是有限指数的反例。某些一般测度甚至不保持零集，因而平移未必能在
几乎处处等价类上良定义；例如点质量测度下在该点相等的两个函数，平移后
在该点未必相等。
\end{solution}

\begin{solution}{translation-discontinuous-linfty}
取 \(f=\mathbf1_{[0,\infty)}\)。对任意 \(h\ne0\)，\(f(\cdot-h)\) 与
\(f\) 在长度 \(|h|\) 的区间上相差 \(1\)，该区间正测。因此
\[
 \|\tau_hf-f\|_\infty=1,
\]
不随 \(h\to0\) 趋零。
\end{solution}

\begin{solution}{mollifier-norm-bound}
对 \(1<p<\infty\)，在质量为1的测度
\(\rho_\varepsilon(y)dy\) 下对 \(|f(x-y)|\cdot1\) 使用 Hölder；
\(p=1\) 时用积分三角不等式，均得
\[
 |(\rho_\varepsilon*f)(x)|^p
 \le\int\rho_\varepsilon(y)|f(x-y)|^p\,dy.
\]
对 \(x\) 积分并用 Tonelli 定理及平移不变性，得到
\[
 \|\rho_\varepsilon*f\|_p^p
 \le\int\rho_\varepsilon(y)\|\tau_yf\|_p^p\,dy
 =\|f\|_p^p.
\]
对 \(p=\infty\)，几乎处处有
\[
 |(\rho_\varepsilon*f)(x)|
 \le\|f\|_\infty\int\rho_\varepsilon(y)\,dy
 =\|f\|_\infty.
\]
\end{solution}

\begin{solution}{mollifier-support}
若 \(x\notin K+\overline{B(0,\varepsilon)}\)，则对每个
\(y\in\operatorname{supp}\rho_\varepsilon\)，有 \(x-y\notin K\)，
故关于 \(y\) 几乎处处有 \(f(x-y)=0\)，因为固定平移及反射保持零集。
因此卷积在该 \(x\) 为零。右端是两个紧集之和，仍为紧集，特别是闭集，
所以非零集的闭包也包含于其中。\(K=\varnothing\) 时 \(f=0\) 几乎处处，
结论同样成立。
\end{solution}

\begin{solution}{density-open-set}
若 \(\Omega=\varnothing\)，空间只有零元，结论成立；若
\(\Omega=\mathbb R^n\)，直接用正文定理。以下假设 \(\Omega\) 非空且
不是全空间。对整数 \(j\ge1\)，取紧集递增穷尽
\[
 K_j=\{x\in\Omega:|x|\le j,\ d(x,\Omega^c)\ge1/j\}.
\]
每个点属于最终的 \(K_j\)，故控制收敛给
\(f\mathbf1_{K_j}\to f\) 于 \(L^p(\Omega)\)。把截断函数零延拓为
\(F_j\in L^p(\mathbb R^n)\)。若 \(K_j=\varnothing\)，取 \(g_j=0\)。
否则 \(d(K_j,\Omega^c)\ge1/j\)，由磨光逼近选
\(0<\varepsilon_j<1/(2j)\)，使
\(\|\rho_{\varepsilon_j}*F_j-F_j\|_p<1/j\)。令
\(g_j=\rho_{\varepsilon_j}*F_j\)。其支撑包含在
\(K_j+\overline{B(0,\varepsilon_j)}\subset\Omega\)，故
\(g_j\in C_c^\infty(\Omega)\)，且
\[
 \|g_j-f\|_{L^p(\Omega)}
 \le 1/j+\|f\mathbf1_{K_j}-f\|_{L^p(\Omega)}\to0.
\]
\end{solution}

\begin{solution}{operator-extension-application}
固定 \(b\in\ell^1(\mathbb Z)\)。先对 \(a\in\ell^p(\mathbb Z)\)
考察有限和
\[
 (S_Na)_n=\sum_{|k|\le N}a_{n-k}b_k.
\]
整数指标的平移保持范数，有限次 Minkowski 给
\[
 \|S_Na\|_p\le\Bigl(\sum_{|k|\le N}|b_k|\Bigr)\|a\|_p.
\]
对 \(1\le p<\infty\)，还有
\[
 \|S_Ma-S_Na\|_p
 \le\|a\|_p\sum_{N<|k|\le M}|b_k|\longrightarrow0\quad(M>N\to\infty).
\]
由 \(\ell^p(\mathbb Z)\) 完备，有限和有范数极限。每个坐标的绝对值
不超过 \(\ell^p\) 范数，所以极限的第 \(n\) 个坐标就是
\(\sum_k a_{n-k}b_k\)。该级数也绝对收敛，因为
\(\|a\|_\infty\le\|a\|_p\)，从而
\(\sum_k|a_{n-k}b_k|\le\|a\|_\infty\|b\|_1\)。令 \(N\to\infty\)，得到
所求 Young 估计。特别地，这给出有限支撑 \(a\) 上的线性算子 \(T_b\)。

截断 \(a^{(j)}=a\mathbf1_{\{|n|\le j\}}\) 使
\(\|a^{(j)}-a\|_p\to0\)，故有限支撑序列稠密。延拓定理给出唯一的
有界延拓。由于
\[
 \left|\sum_k(a^{(j)}_{n-k}-a_{n-k})b_k\right|
 \le\|b\|_1\|a^{(j)}-a\|_\infty
 \le\|b\|_1\|a^{(j)}-a\|_p\to0,
\]
这个延拓与上述逐坐标卷积一致。

当 \(p=\infty\)，同一绝对收敛估计直接给出
\(\|T_ba\|_\infty\le\|b\|_1\|a\|_\infty\) 和线性性，不用稠密延拓。
常数序列1与任何有限支撑序列的距离至少为1，所以后者在
\(\ell^\infty(\mathbb Z)\) 中不稠密。
\end{solution}
